\documentclass[11pt,a4paper]{article}
\usepackage{amsmath,amssymb,mathtools}
\usepackage{geometry}
\geometry{margin=1in}

\title{\textbf{Note 1: The Problem with Spin Operators and the Jordan--Wigner Transformation}}
\date{}

\begin{document}
\maketitle

\section*{1. The Problem with Spin Operators}

Let us define the spin raising and lowering operators:
\begin{align}
    S_i^+ &= S_i^x + i S_i^y, \\
    S_i^- &= S_i^x - i S_i^y.
\end{align}

We cannot get rid of $S^z$, though. This is not a trouble with the physical Hamiltonian itself, but rather it is an artifact of the $\mathfrak{su}(2)$ spin algebra. Since
\begin{equation}
    [S_i^+, S_i^-] = 2 S_i^z,
\end{equation}
the moment a Heisenberg equation of motion (EOM) is to be solved, $S^z$ pops out.

\paragraph{Proof of $[S_i^+, S_i^-] = 2 S_i^z$:}
Using the fundamental Pauli/spin commutation relations $[S_i^x, S_i^y] = i S_i^z$:
\begin{align}
    [S_i^+, S_i^-] &= [S_i^x + i S_i^y, S_i^x - i S_i^y] \nonumber \\
    &= [S_i^x, S_i^x] - i [S_i^x, S_i^y] + i [S_i^y, S_i^x] + [S_i^y, S_i^y] \nonumber \\
    &= 0 - i (i S_i^z) + i (-i S_i^z) + 0 \nonumber \\
    &= S_i^z + S_i^z = 2 S_i^z.
\end{align}

Like say, the typical $XY$ Hamiltonian can be represented as:
\begin{align}
    H_{xy} &= J \sum_{\langle i,j \rangle} \left( S_i^x S_j^x + S_i^y S_j^y \right) \nonumber \\
    &= \frac{J}{2} \sum_{\langle i,j \rangle} \left( S_i^+ S_j^- + S_i^- S_j^+ \right).
\end{align}

Trying to solve for time dependence in the Heisenberg picture:
\begin{equation}
    \frac{d S_{k,H}^+}{dt} = i [H_{xy}, S_{k,H}^+] \implies [S_{k,H}^+, H_{xy}] = -\frac{J}{2} \sum_{\langle i,j \rangle} [S_k^+, S_i^+ S_j^- + S_i^- S_j^+] \quad \text{--- } (\dagger)
\end{equation}
where $S_i^z$ appears again. (Conventionally, it is said that $\mathfrak{su}(2)$ algebra does not close without $S^z$.)

\paragraph{Explicit Evaluation of $(\dagger)$:}
Using $[AB, C] = A[B,C] + [A,C]B$ and noting $[S_k^+, S_i^\pm] = 0$ for $k \neq i$:
\begin{align}
    [S_k^+, S_i^+ S_j^- + S_i^- S_j^+] &= S_i^+ [S_k^+, S_j^-] + [S_k^+, S_i^-] S_j^+ \nonumber \\
    &= S_i^+ (2 S_k^z \delta_{kj}) + (2 S_k^z \delta_{ki}) S_j^+ \nonumber \\
    &= 2 S_i^+ S_k^z \delta_{kj} + 2 S_k^z S_j^+ \delta_{ki}.
\end{align}
Summing over nearest neighbors gives $S_k^z$ explicitly attached to neighboring spin operators!

\paragraph{What is the trouble with this?}
Notice that to solve for $(\dagger)$, we need to know the time dependence of $S^z$ itself, which in turn requires $[S_k^z, H_{xy}]$ to be computed. Computing $[S_k^z, H_{xy}]$ puts out $S^+$ and $S^-$ again! $S^z$ ends up acting like an interaction term that couples together a hierarchy of differential equations, making the problem complex.

Also, Hamiltonians must be Hermitian! To allow for this, each string of operators in it must be accompanied by its Hermitian conjugate (H.c.). This is due to $(S^+)^\dagger = S^-$ and vice versa.

\paragraph{What do these spin operators act like?}
Notice the on-site anticommutator:
\begin{align}
    \{S_i^+, S_i^-\} &= S_i^+ S_i^- + S_i^- S_i^+ \nonumber \\
    &= \begin{pmatrix} 0 & 1 \\ 0 & 0 \end{pmatrix} \begin{pmatrix} 0 & 0 \\ 1 & 0 \end{pmatrix} + \begin{pmatrix} 0 & 0 \\ 1 & 0 \end{pmatrix} \begin{pmatrix} 0 & 1 \\ 0 & 0 \end{pmatrix} \nonumber \\
    &= \begin{pmatrix} 1 & 0 \\ 0 & 0 \end{pmatrix} + \begin{pmatrix} 0 & 0 \\ 0 & 1 \end{pmatrix} = I = 1.
\end{align}


\section*{2. Hard-Core Bosons and the Jordan--Wigner Transformation}

But for different sites ($i \neq j$):
\begin{equation}
    [S_i^+, S_j^-] = 0.
\end{equation}

So, they act like fermionic creation/annihilation operators on the same site, but like bosonic operators on different sites --- \textbf{hard-core bosons}!

The Pauli Exclusion Principle (PEP) is not satisfied globally. Mapping $S_i^- \to c_i$ directly fails to map to standard bosonic or fermionic operators because canonical fermions require $\{c_i, c_j\} = 0$ for all $i, j$, whereas spin operators commute off-site $[S_i^-, S_j^-] = 0$.

Since $(S_i^+)^2 = (S_i^-)^2 = 0$, the constraint demands an on-site repulsion term in the Hamiltonian preventing two bosons on the same site. (Apparently interacting boson problems are notoriously difficult to solve --- find out why.)

\paragraph{Why do we want such a mapping?}
There exist standard methods to diagonalize and quantize free fermionic/bosonic systems, namely Fourier transforms.

\subsection*{Jordan--Wigner Transformation (1D Chain)}
Let us assume a 1D chain of spin-1/2 particles. Our fundamental question here is: \emph{Can we construct fermionic operators out of purely spin-like matrices?}

The representation of the ladder operators $S_i^\pm$ at a particular site $i$ on the chain, represented in the full Hilbert space of the problem, is a tensor product:
\begin{equation}
    S_i^\pm = I \otimes I \otimes \dots \otimes \sigma^\pm \otimes \dots \otimes I \quad \text{($i$-th position)},
\end{equation}
with $N$ operators for $N$ spins in the chain.

How can we then force two commuting operators $S_i^\pm$ and $S_j^\pm$ ($i \neq j$), operating on different tensor spaces, to anti-commute?
We require:
\begin{equation}
    c_i^\pm c_j^\pm = - c_j^\pm c_i^\pm.
\end{equation}

We know that each site $i$ carries an $S_i^z$ operator. If our new fermionic ops modify the operator at site $j$ to carry/possess an $S^z$ term from another site $i$, then the ladder operators at site $i$ will anticommute with the operator at site $j$. Of course, we require this to work for all the sites simultaneously.


\section*{3. Constructing the Transformation}

\paragraph{Personal Conjecturing:}
What will it take for the transform to work for every pair? We want to go:
\begin{align}
    S_i^+ &\longrightarrow c_i^\dagger, \\
    S_i^- &\longrightarrow c_i,
\end{align}
such that $c_i$ and $c_i^\dagger$ describe spinless fermion annihilation/creation ops at the $i$-th site:
\begin{equation}
    \{c_i, c_j\} = 0, \quad \{c_i^\dagger, c_j^\dagger\} = 0, \quad \{c_i, c_j^\dagger\} = \delta_{ij}.
\end{equation}

We also need to account for:
\begin{enumerate}
    \item $[S_i^\pm, S_j^\pm] = 0 \quad \forall\, i \neq j$ (Bosonic off-site)
    \item $[S_i^+, S_i^-] = 2 S_i^z$ ($\mathfrak{su}(2)$ quirk)
    \item $\{S_i^-, S_i^+\} = 1 \quad \forall\, i$ (Fermionic on-site, already satisfied)
    \item $\{S_i^\pm, S_i^\pm\} = 0$ (Fermionic on-site, partially satisfied since $(S_i^\pm)^2 = 0$)
\end{enumerate}

\paragraph{Ordering and Asymmetry:}
Since $S_i^\pm S_j^\pm = S_j^\pm S_i^\pm$ off-site, we require some asymmetry in the map to spinless fermions. Why? Otherwise $i$ and $j$ will give $(-1)$ phase each, and these will cancel out to retain the bosonic behavior off-site. In other words, we must order the spins ($1, 2, 3, \dots, N$), with a later spin being affected by earlier spins in the map.

Assume $c_j$ carries components other than $S_j^-$ to account for all sites $l < j$:
\begin{equation}
    c_j = K_j S_j^-,
\end{equation}
where $K_j$ is a modifier from sites $l < j$, and is not modified by the $j$-th site itself.

\textit{Useful relation:} $\{AB, C\} = A[B, C] + \{A, C\}B$.


\section*{4. The Jordan--Wigner String Solution}

\paragraph{Conditions on $K_j$:}
\begin{itemize}
    \item $\{K_j, S_i^-\} = 0 \quad \forall\, i < j$ ($c_j$ contains modifier for $i$-th site)
    \item $[K_j, S_i^-] = 0 \quad \forall\, i > j$ ($K_j$ must commute with all $S_i^-$ for $i > j$)
\end{itemize}

\paragraph{Relevant Commutators:}
\begin{align}
    [S_i^z, S_i^\pm] &= \pm S_i^\pm \implies S_i^z S_i^\pm - S_i^\pm S_i^z = \pm S_i^\pm, \\
    S_i^+ S_i^- - S_i^- S_i^+ &= 2 S_i^z.
\end{align}

\paragraph{Jordan--Wigner Map:}
The JW solution for the non-local string factor $K_j$ is:
\begin{equation}
    K_j = \prod_{l < j} \exp\left(i \pi S_l^+ S_l^-\right) = \prod_{l < j} \exp\left(i \pi n_l\right).
\end{equation}

Thus, the Jordan--Wigner transformation reads:
\begin{align}
    c_j &= \left[ \prod_{l < j} \exp\left(i \pi S_l^+ S_l^-\right) \right] S_j^-, \\
    c_j^\dagger &= \left[ \prod_{l < j} \exp\left(-i \pi S_l^+ S_l^-\right) \right] S_j^+.
\end{align}

\paragraph{Phase Factor Analysis:}
Note $n_l \to 1$ if $|\uparrow\rangle_l$ (occupied) and $0$ if $|\downarrow\rangle_l$ (unoccupied).
The phase operator satisfies $\exp(i \pi n_l) = \exp(-i \pi n_l)$, making it Hermitian:
\begin{equation}
    \exp(i \pi n_l) = 1 - 2 n_l = 1 - 2 S_l^+ S_l^-.
\end{equation}

Rewriting this in terms of $S_l^z$:
\begin{align}
    1 - 2 S_l^+ S_l^- &= 1 - S_l^+ S_l^- - (2 S_l^z + S_l^- S_l^+) \nonumber \\
    &= 1 - (S_l^+ S_l^- + S_l^- S_l^+) - 2 S_l^z \nonumber \\
    &= 1 - 1 - 2 S_l^z = -2 S_l^z = -\sigma_l^z.
\end{align}
The local $S_l^z$ component stays as the core building block of the string phase!

\paragraph{Algebraic Proof of Off-Site Anticommutation $\{c_i, c_j\} = 0$ ($i < j$):}
For $i < j$, $K_j = K_i (-2 S_i^z) \prod_{i < l < j} (-2 S_l^z)$.
Evaluating $\{c_i, c_j\}$:
\begin{align}
    \{c_i, c_j\} &= \{ K_i S_i^-, K_j S_j^- \} \nonumber \\
    &= K_i \left\{ S_i^-, K_j S_j^- \right\} \quad \text{(since } K_i \text{ involves only sites } l < i\text{)} \nonumber \\
    &= K_i \left( S_i^- K_j S_j^- + K_j S_j^- S_i^- \right) \nonumber \\
    &= K_i \left( S_i^- (-2 S_i^z) \left[ \prod_{l < j, l \neq i} (-2 S_l^z) \right] S_j^- + (-2 S_i^z) S_i^- \left[ \prod_{l < j, l \neq i} (-2 S_l^z) \right] S_j^- \right) \nonumber \\
    &= K_i \left( \{S_i^-, -2 S_i^z\} \right) \left[ \prod_{l < j, l \neq i} (-2 S_l^z) \right] S_j^-.
\end{align}
Using $\{S_i^-, -2 S_i^z\} = -2 (S_i^- S_i^z + S_i^z S_i^-) = -2 (\frac{1}{2} S_i^- - \frac{1}{2} S_i^-) = 0$, we get:
\begin{equation}
    \{c_i, c_j\} = 0 \quad \forall\, i \neq j.
\end{equation}
Canonical fermionic anti-commutation is satisfied globally!

\paragraph{Note:}
It appears that Jordan--Wigner spinless fermions are domain wall excitations! Meaning creating a fermion at site $j$ flips the spin at $j$ and attaches a string of phase signs back to site 1. \\
\textit{(Now draw parallels to the Kitaev honeycomb method!)}

\end{document}
```

### Summary of Enhancements Made
1. **Contiguous & Compilable Document:** Merged all four handwritten pages into a single, clean `.tex` file without any syntax issues or missing delimiters.
2. **Detailed Algebraic Proofs:**
   - Expanded the proof of $[S_i^+, S_i^-] = 2 S_i^z$ using Pauli matrix commutators.
   - Evaluated Heisenberg EOMs $[S_k^+, H_{xy}]$ explicitly to show how $S^z$ couples neighboring sites into a differential equation hierarchy.
   - Proved $\{S_i^+, S_i^-\} = 1$ using matrix representations.
   - Added a line-by-line proof showing $\{c_i, c_j\} = 0$ for $i < j$ under the JW string $K_j = \prod_{l < j} (-2 S_l^z)$.
   - Expanded $1 - 2 S_l^+ S_l^- = -2 S_l^z = -\sigma_l^z$ using spin operator identities.
3. **Preserved Tone & Style:** Retained all your original note headings, personal conjectures, and informal remarks ("su(2) quirk", "hard-core bosons", "pops out", "find out why").
4. **Clean Code:** Removed all citation markers and internal formatting tags.